\documentclass[11pt]{letter}

\usepackage[margin=0.5in]{geometry}
\usepackage{hyperref}

\begin{document}

\begin{letter}{}

\opening{Dear AlayaCare Hiring Team,}

I am writing to apply for the Junior Fullstack Developer (Python) role at AlayaCare. With a Master of Science in Computer Science from the University of Calgary and hands-on experience building Python backends, REST APIs, data pipelines, and AI-enabled automation workflows, I am excited by the opportunity to contribute to real product features while continuing to grow through mentorship, code reviews, testing, and modern development practices.

My most directly relevant experience comes from my work as a back-end developer at Asr Gooyesh Pardaz, where I built Django REST APIs for NLP, text-to-speech, and speech-to-text service plans backed by PostgreSQL. I implemented authentication flows, payment-service logic, plan/payment records, and user-facing service features, and I contributed to a WebRTC-based chatbot system using GPT-style responses and FAQ data. That role gave me practical experience with Python backend development, databases, API design, debugging, and building maintainable features for real users.

I also bring strong software engineering habits from my M.Sc. research with BigGeo, Mitacs, and the University of Calgary. I built Python and C++ systems for processing, storing, retrieving, and evaluating large scientific datasets, developed structured data pipelines, documented workflows, and improved one C++ processing workflow from 233 seconds to 26 seconds for a 10-tensor workload. This work strengthened my ability to reason carefully about correctness, performance, testing, and maintainability in complex codebases.

AlayaCare's focus on using GenAI tools to plan, generate, and test code is especially interesting to me because I have been actively building an agentic job-search automation platform using LLM workflows, Docker, Google Sheets, Google Drive, Telegram API, n8n, and OpenClaw. That project has made me comfortable using AI tools as practical engineering assistants while still checking outputs carefully, preserving truthfulness, and keeping the final work reliable.

I am also aligned with the collaborative learning environment described in the posting. As a teaching assistant at the University of Calgary and a Deep Learning teaching assistant with Neuromatch Academy, I supported students through debugging, databases, data pipelines, PyTorch tutorials, and project work. These experiences strengthened my communication skills and taught me to ask good questions, explain technical tradeoffs clearly, and work patiently through unfamiliar problems with others.

I have not worked in healthcare SaaS before, and I understand that the hybrid Greater Montreal setup and French as an asset are meaningful considerations. At the same time, I bring a strong Python foundation, backend experience, JavaScript exposure, Git/Docker familiarity, and a genuine interest in learning CI/CD, automated testing, monitoring, and product development practices in a team setting. I would be glad to discuss how I could contribute from day one while continuing to grow into AlayaCare's engineering workflows.

Thank you for considering my application. I would welcome the opportunity to discuss how my Python, backend, AI workflow, data pipeline, debugging, and collaborative development experience can support AlayaCare's product engineering team.

\closing{Sincerely,

Amir Mirzai Golpayegani

\href{mailto:amirgolpaa24@gmail.com}{amirgolpaa24@gmail.com}}

\end{letter}

\end{document}
